\paragraph{Pointwise bounds under a core operation.}
Even when the upper bound is itself a core operation, separate bounded-profile
tests need not have a common witness. A single positive profile bit already
fails to glue on three variables. The failure also occurs with exact profiles
and persists for the value bounds required by a single-owner occurrence type.

\begin{lemma}[Singleton-profile absorption]\label{core:lem:singleton-profile-absorption}
For a core operation $H$, an input $j$, and an owner $r$,
\[
 \Phi_r(H)(\{j\})=1
 \quad\Longleftrightarrow\quad H=j\star_r H.
\]
Equivalently, at every valuation $w$, if $w_j=\{r\}$ then
$H(w)=\{r\}$; otherwise $H(w)=\{r\}$ or $w_j\subseteq H(w)$.
\end{lemma}
\begin{proof}
For the forward direction, induct on a core term with the positive profile
bit. A projection must be $j$. At an $r$-root, one child has the same
positive bit. If either child is the $r$ singleton, so is the root;
otherwise the inductive child contains $w_j$, and union retains it.
At a foreign $s$-root, both children have the bit. If both are the $r$
singleton, so is the root. Otherwise one contains $w_j$. An $s$-absorbing
child also contains $w_j$, so absorption can occur only with
$w_j=\{s\}$; without absorption union retains $w_j$.
In addition, $w_j=\{r\}$ forces the $r$ singleton by profile monotonicity.
These are precisely the pointwise conditions for $H=j\star_r H$.
Conversely that identity forces $H=\{r\}$ whenever $j=\{r\}$, proving
the positive profile bit.
\end{proof}

\begin{proposition}[A positive singleton-profile bit does not glue]
\label{core:prop:singleton-profile-bound-gluing}
Fix distinct owners $a,b,c$, and put
\[
 C=(x\star_b j)\star_a(j\star_c y),\qquad
 H_0=j,\qquad H_1=j\star_b x.
\]
Both helpers satisfy $\Phi_b(H_i)(\{j\})=1$, and at every valuation $w$
at least one satisfies $H_i(w)\subseteq C(w)$. Nevertheless no core
operation $H\subseteq C$ satisfies $\Phi_b(H)(\{j\})=1$.
\end{proposition}
\begin{proof}
The profile bits are immediate. If $j\subseteq C$, use $H_0$.
Otherwise the outer $a$-gate cannot absorb: its $a$-profile is
$j(x\lor y)$, so an $a$ output would contain $j=\{a\}$.
Both children must therefore omit some member of $j$, forcing their
respective $b$- and $c$-absorptions. Thus $C=\{b,c\}$ and
$j\ne\{b\},\{c\}$, whence $x=\{b\}$ and $y=\{c\}$.
Now $H_1=\{b\}\subseteq C$.

Suppose a bounded $H$ has the positive bit. At the atomic input
$(x,j,y)=(\{b\},\{a\},\{c\})$, the bound is $\{b,c\}$.
The absorption lemma forces $H=\{b\}$, giving
$\Phi_b(H)(\{x\})=1$. At
$(x,j,y)=(\{b\},\{a\},\{a\})$, the bound is $\{a\}$, giving
$\Phi_a(H)(\{j,y\})=1$. These positive coalitions are disjoint,
contradicting the common-source property.
\end{proof}

We next keep the helpers' entire profiles fixed.

Fix distinct owners $a,b,c$, and abbreviate
\[
 P=x\star_a y,\quad S=x\star_b y,\quad
 Q=P\star_c z,\quad R=z\star_c x.
\]
Put
\begin{equation}\label{core:eq:profile-gluing-words}
 C=(S\star_b z)\star_a Q,\qquad
 H_0=(R\star_b S)\star_a Q,\qquad
 H_1=(R\star_a P)\star_c(P\star_b z).
\end{equation}
The profile notation below uses $x,y,z$ for coalition membership bits.

\begin{proposition}[Pointwise profile bounds do not glue]
\label{core:prop:profile-bound-gluing}
For every $K\geq3$, the admissible profile family
\[
 \phi_a=z(x\lor y),\qquad \phi_c=xy,\qquad
 \phi_r=xyz\quad(r\notin\{a,c\})
\]
contains all positive profile bits of $C$. At every valuation $w$, one of
$H_0,H_1$ has value contained in $C(w)$, and both have exact profiles
$\phi$. Nevertheless no core operation $H\subseteq C$ has these profiles.
In fact no such $H$ has $\Phi_c(H)=xy$.
\end{proposition}
\begin{proof}
Profile recursion gives the stated profiles of both helpers. The profiles
of $C$ are $\Phi_a(C)=z(x\lor y)$ and $\Phi_r(C)=xyz$ for $r\ne a$.

We prove the pointwise cover for arbitrary nonempty menus, allowing extra
owners. If $z=\{c\}$ and either $x=\{b\}$ or $y=\{b\}$, then
$C=H_0=\{b,c\}$. Otherwise $H_1\subseteq C$ as follows.
If $P\ne\{a\}$, then $P=x\cup y$; neither $S\star_b z$ nor $Q$ is the
$a$ singleton, so $C=(S\star_b z)\cup Q$. Also $R\ne\{a\}$ and
$H_1\subseteq P\cup z$. One of $S\star_b z,Q$ equals $P\cup z$ unless
both the $b$ and $c$ absorptions occur. Outside the excluded case, simultaneous
absorption requires $z=\{b\}$ and $P=\{c\}$, where
$C=P\cup z=\{b,c\}$ anyway. If $P=\{a\}$, direct evaluation gives
$H_1=\{a\}\cup z$. For $z=\{a\},\{b\}$ this equals $C$.
For $z=\{c\}$ outside the excluded case, $S\star_b z$ contains
$\{a,c\}$. For every other $z$, the ambiguous value
$Q=\{a\}\cup z$ is contained in $C$.

For nonexistence, we establish the following two-row implication for every
core term $t$:
\begin{equation}\label{core:eq:two-row-profile-implication}
 \begin{gathered}
 t(\{a\},\{b\},\{c\})\subseteq\{b,c\},\qquad
 t(\{a\},\{c\},\{b\})\subseteq\{a,b\}
 \end{gathered}
 \quad\Longrightarrow\quad
 \Phi_c(t)(\{x,y\})\leq\Phi_c(t)(\{z\}).
\end{equation}
Call these inputs $u,v$. Suppose a counterexample exists, and descend its
tree retaining both bounds, the positive $c$-bit on $\{x,y\}$, and the
negative $c$-bit on $\{z\}$. Its value at $v$ is ambiguous: it cannot
be the $c$ singleton by the bound, the $a$ singleton because the
$a$-atomic coalition $\{x\}$ is the same at $u$, or the $b$ singleton
because its coalition $\{z\}$ is disjoint from the positive $c$-coalition
$\{x,y\}$. Thus both children satisfy the bound at $v$.
At $u$ a singleton absorber can only be $b$: the bound excludes $a$,
and the negative $c$-bit excludes $c$. If a $b$-root has a $b$-absorbing
child at $u$, choose it. It retains the positive $c$-bit, while its
positive $b$-bit on $\{y\}$ excludes the $c$-bit on $\{z\}$ by common
source. In every other case both bounds pass to both children; the
profile OR/AND recurrences select a child retaining the positive and
negative $c$-bits. Extra-owner gates act as union at these two inputs.
Descent reaches a projection, but the only projection satisfying both
bounds is $z$, whose $c$-bit on $\{x,y\}$ is false. This contradiction
proves \eqref{core:eq:two-row-profile-implication}.

Finally $C(u)=\{b,c\}$ and $C(v)=\{a,b\}$. The profile $xy$ violates
the implication, excluding every $H\subseteq C$ with that profile.
\end{proof}
